\documentclass[11pt]{article}
% Shared preamble of the RMC kernel write-ups (\input{preamble} after \documentclass).
\usepackage[margin=1in]{geometry}
\usepackage{amsmath,amssymb,bm}
\usepackage{graphicx}
\usepackage{booktabs}
\usepackage{hyperref}

\newcommand{\dd}{\,\mathrm{d}}
\newcommand{\Ein}{E}
\newcommand{\Eout}{E'}
\newcommand{\Ecm}{E_{c}}
\newcommand{\mucm}{\mu_{c}}
\newcommand{\mulab}{\mu_0}

\newcommand{\acespec}[1]{J.~L.~Conlin and P.~Romano,
\href{https://mcnp-green.lanl.gov/pdf_files/TechReport_2019_LANL_LA-UR-19-29016_ConlinRomano.pdf}
{\emph{A Compact ENDF (ACE) Format Specification}}, LA-UR-19-29016 (2019),
Sec.~#1}
\newcommand{\ENDFspec}[1]{\href{https://www.nndc.bnl.gov/endfdocs/ENDF-102-2023.pdf}
{\emph{ENDF-6 Formats Manual}} (2023), #1}

% Common closing section: how the verification figures are produced
\newcommand{\verificationmethod}{%
Each figure is produced by \texttt{verification/verify\_laws.py}. For an incident
energy $\Ein$ along $+z$, $10^6$ collisions are sampled with MC/DC's own reaction
sampler (\texttt{sample\_elastic\_scattering}, \texttt{sample\_inelastic\_scattering}
or \texttt{sample\_fission}), and every emitted neutron is histogrammed in
$48 \times 24$ bins of $(\Eout, \mulab)$, with $\mulab$ the lab cosine to $+z$. The
inverted kernel used by RMC is integrated over the same bins (Gauss--Legendre panels;
two-body lines are split where $\mulab^*(\Eout)$ crosses a cosine edge), multiplied by
the number of collisions, and compared with a Pearson-type $\chi^2$ diagnostic
over bins with at least 5 expected counts (the rest pooled into one bin).
The existing tables list the script's $\chi^2/\mathrm{dof}$ and nominal $p$-value,
the emitted neutrons per collision (sampled vs.\ the integral of the
inverted kernel), and the largest relative change of a bin integral when the
quadrature panels are doubled. Left panels: outgoing-energy marginal; middle:
lab-cosine marginal; right: per-bin $z = (\text{sampled} - \text{inverted}) /
\sqrt{\text{inverted}}$.

\noindent These figures check agreement with the MC/DC sampler, not independent
physical validity. Their existing $p$-values are not calibrated goodness-of-fit
probabilities: the script uses $B$ degrees of freedom, whereas a fixed-total
multinomial histogram with $B$ categories and no fitted parameters uses $B-1$.
Random multiplicities and spectrum assignments can also change the collision-level
covariance from independent Poisson counts; $z$ is then only a count diagnostic.
See \href{https://docs.scipy.org/doc/scipy/reference/generated/scipy.stats.chisquare.html}
{the Pearson test's sampling assumptions}.
Energy bins use the sample extrema, so unsampled tails are excluded and the
reported inverted yield is an integral over that observed window. A formal test
needs independently specified bins, explicit tail categories, and a covariance or
resampling model appropriate to independent collisions. Zero expected density
with observed emissions must be reported rather than discarded. Doubling quadrature
panels gives a convergence indicator, not a rigorous error bound. The stored
figures and tables have not been regenerated for these qualifications.}

\usepackage{amsthm}

\newtheorem{proposition}{Proposition}
\newcommand{\tpsi}{\tilde\psi}
\newcommand{\teps}{\tilde\epsilon}
\newcommand{\Sbar}{\bar S}
\newcommand{\Sigt}{\Sigma_t}
\newcommand{\Om}{\boldsymbol\Omega}
\newcommand{\That}{\hat T}
\newcommand{\E}{\mathbb{E}}

\title{Full-residual correction passes in 3D with a tabulated in-scatter kernel}
\date{}

\begin{document}
\maketitle

\noindent This note gives the mathematics of the correction passes of the 3D
Cartesian solver (\texttt{mcdc/rmc/correction3d.py}, \texttt{driver3d.run3d} with
\texttt{N\_correction}). Notation and the trial space are those of
\texttt{cartesian\_3d/cartesian\_3d.tex}: cells $V_k$, energy bins $g$, angular
patches $A_j$ ($\mathrm d\Om = \mathrm d\mu\,\mathrm d\varphi$), and $\tpsi$ constant
on each $(k,g,j)$.

\section{Why a correction pass, and why it is unbiased}

The identities in this section assume the exact kernel $T$. The implementation
uses a tabulated interpolant $\hat T$ and retains the deterministic bias described
under Interpolation below; matching bin averages does not remove it.

The collision-only iterations sample the residual with the binned in-scatter
$\Sbar = M\tpsi/(\Delta E_g|A_j|)$ in place of the in-scatter $T[\tpsi](E,\Om)$ of the
trial function. Their fixed point $\tpsi^*$ therefore carries the in-bin shape bias
of the constant trial space. A correction pass samples the full residual
\begin{equation}
  r = q - L\tpsi = \underbrace{Q + T[\tpsi] - \Sigt(E)\,\tpsi}_{\text{inside each cell}}
  \; + \; r_e ,
\end{equation}
with $r_e$ the face term of \texttt{cartesian\_3d.tex}, transports it with MC/DC's
unchanged operator $L$, and adds the projected response:
\begin{equation}
  \tpsi + \Pi L^{-1}(q - L\tpsi) = \tpsi + \Pi\psi - \Pi\tpsi = \Pi\psi
  \label{eq:unbiased}
\end{equation}
for any $\tpsi$ in the trial space ($\Pi\tpsi=\tpsi$), where $\psi=L^{-1}q$. So the
expectation of one pass, given $\tpsi$, is the exact bin average $\Pi\psi$; this is
the RMC update of [1--3]. Here $T$ is the scattering \emph{and} prompt-fission
emission: fission born inside the window (the fast window) is part of $T$, and its
binned part is part of $\Sbar$ through the transfer moments. In the thermal window
fission counts as capture.

\paragraph{Error estimate.} By the law of total variance [4, Ch.~8],
$\operatorname{Var}(\tpsi+\teps) = \E[\operatorname{Var}(\teps\mid\tpsi)] +
\operatorname{Var}(\E[\tpsi+\teps\mid\tpsi])$, and the second term vanishes by
\eqref{eq:unbiased}. The history standard error of the pass's tally, which estimates
$\operatorname{Var}(\teps\mid\tpsi)$, is therefore the error of the corrected result
(\texttt{RMC3DResult.psi\_sdev}). Several passes from the same $\tpsi^*$ are averaged,
as in the slab solver.
For the implemented interpolant, the conditional corrected mean includes a
trial-dependent interpolation error. The history standard error is conditional
sampling uncertainty, not total error or an interpolation-error estimate.

\section{The in-scatter of the trial function}

Inside cell $k$ of material $m$, $\tpsi$ is constant on each incident bin, so
\begin{equation}
  T(E,\Om) = \sum_{g',j'} \tpsi_{kg'j'}\, U^m_{g'j'}(E,\Om),\qquad
  U^m_{g'j'}(E,\Om) = \int_{g'}\!\mathrm dE' \sum_r N_r\sigma_r(E')
     \int_{A_{j'}}\!\mathrm d\Om'\, \frac{f_r(E'\!\to E,\Om'\cdot\Om)}{2\pi},
  \label{eq:U}
\end{equation}
with $f_r$ the lab emission density of reaction $r$ per unit outgoing energy and lab
cosine $\mu_0$ (yield included, \texttt{mcdc.rmc.reaction}). The kernel depends on
the directions only through $\mu_0=\Om'\cdot\Om$ because the medium is isotropic
[5, Ch.~1].

\begin{proposition}[Patch integral of a rotationally symmetric kernel]
For any integrable $h(\mu_0)$,
\begin{equation}
  \int_{A_{j'}} \frac{h(\Om'\cdot\Om)}{2\pi}\,\mathrm d\Om' =
  \int_{-1}^{1} h(\mu_0)\, F_{j'}(\Om,\mu_0)\,\mathrm d\mu_0 ,
\end{equation}
with $F_{j'}(\Om,\mu_0)$ the fraction of the circle
$\{\Om' : \Om'\cdot\Om = \mu_0\}$ that lies in $A_{j'}$.
\end{proposition}
\begin{proof}
In polar coordinates about $\Om$,
$\Om' = \mu_0\Om + \sqrt{1-\mu_0^2}\,(\cos\omega\,\mathbf e_1 + \sin\omega\,\mathbf e_2)$
and $\mathrm d\Om' = \mathrm d\mu_0\,\mathrm d\omega$; at fixed $\mu_0$ the
$\omega$-measure of $\{\Om'\in A_{j'}\}$ is $2\pi F_{j'}(\Om,\mu_0)$.
\end{proof}

\noindent $F_{j'}$ is computed exactly from the circle's crossings of the patch
edges (\texttt{patches.arc\_fractions}, \texttt{cartesian\_3d.tex}), so no azimuthal
quadrature appears. Three cases follow.
\begin{itemize}
  \item \emph{Delta laws} (elastic on a target at rest, level scattering):
  $f = g(E\mid E')\,\delta(\mu_0-\mu_0^*(E',E))$, so the inner integral is
  $g(E\mid E')\,F_{j'}(\Om,\mu_0^*)$, one arc-fraction evaluation.
  \item \emph{Isotropic lab emission} (fission): $f(E'\!\to E,\mu_0)$ does not
  depend on $\mu_0$, and setting $h=1$ in the proposition gives
  $\int_{-1}^1F_{j'}\,\mathrm d\mu_0 = |A_{j'}|/2\pi$, so the inner integral is
  exactly $f\,|A_{j'}|/2\pi$.
  \item \emph{Continuous laws} (continuum inelastic, free gas): 16-point
  Gauss--Legendre [6] over the supported $\mu_0$ range (for centre-of-mass laws, the
  range where the centre-of-mass energy lies in the spectrum's support), with the
  cosine map $\mu_0 = c + (d-c)(1-\cos\pi t)/2$ that absorbs square-root endpoint
  behaviour, as in the slab kernels.
\end{itemize}

\paragraph{Consistency with the transfer moments.} Integrating \eqref{eq:U} over
$E\in g$ and $\Om\in A_j$ and using
$\int_{A_j}F_{j'}(\Om,\mu_0)\,\mathrm d\Om = A_{j'j}(\mu_0) = A_{jj'}(\mu_0)$
(\texttt{cartesian\_3d.tex}) gives the transfer moments: $\int_g\!\int_{A_j} T =
(M\tpsi)_{gj} = \Delta E_g|A_j|\,\Sbar_{gj}$. $\Sbar$ is the exact bin average of $T$.

\section{Tabulation}

$U^m$ depends only on the nuclear data and the grids, so it is computed once per
material and cached on disk (\texttt{correction-table-<hash>.h5}, keyed by the data
contents, the grids and the nodes).

\paragraph{Symmetry.} Let $\sigma$ be an orthogonal map of the sphere that maps the
patch layout onto itself: the rotation about $\hat{\mathbf z}$ by the azimuthal patch
width (uniform $\varphi$ edges) and the reflection $\mu\to-\mu$ (symmetric $\mu$
edges), and their products. Then $F_{\sigma j'}(\sigma\Om,\mu_0) = F_{j'}(\Om,\mu_0)$,
so $U_{\sigma j'}(E,\sigma\Om) = U_{j'}(E,\Om)$ and
\begin{equation}
  T(E,\sigma\Om) = \sum_{g',j''} \tpsi_{kg'\,\sigma(j'')}\, U_{g'j''}(E,\Om).
\end{equation}
Every patch is the image $\sigma_j(f_j)$ of a \emph{fundamental} patch $f_j$, so $U$
is tabulated only for directions in the fundamental patches, and $T$ on patch $j$ is
the table applied to $\tpsi$ permuted by $\sigma_j$. For octant patches the eight
patches are images of one, and the table is eight times smaller.

\paragraph{Incident-energy integral.} For each node energy $E$, reaction and incident
bin $g'$, the $E'$ integral uses adaptive Gauss--Kronrod (G7--K15) [7, 8] on the
summed emission, to a relative tolerance of $10^{-6}$, from four initial pieces;
the per-direction contributions are accumulated with the accepted Kronrod weights.
For elastic scattering on a target at rest, outgoing energies lie in $[\alpha E',E']$
with $\alpha = ((A-1)/(A+1))^2$ [9, Ch.~2], so the integral is clipped to
$E'\in[E, E/\alpha]$; other reactions skip incident bins whose kinematic support
(\texttt{reaction.lab\_energy\_support}) cannot reach $E$.

\section{Interpolation}

\paragraph{Nodes.} In energy, each bin has its own nodes: both edges, $n_E$ uniform
intervals, and the points $\alpha_n E_{g''}$ for every nuclide $n$ and every bin edge
$E_{g''}$, where an incident bin edge enters the elastic range $[E,E/\alpha]$ and $T$
has a kink. (Nodes closer than 2\% of the bin width are merged.) In direction, each
patch carries an $n_\mu\times n_\varphi$ grid including its edges.

\paragraph{Interpolant and shift.} On the sub-cells of this node grid, $T$ is
interpolated trilinearly in $(E,\mu,\varphi)$. Piecewise-linear interpolation of a
smooth function has error at most $h^2\max|f''|/8$ per direction [10, Ch.~3], so the
error is second order away from kinks. Because $\mathrm d\Om = \mathrm d\mu\,
\mathrm d\varphi$, the bin average of the interpolant is the product of trapezoid
rules and is exact. The interpolant $\That$ is shifted by a constant per (cell, bin,
patch) so that its bin average equals $\Sbar$:
\begin{equation}
  \That = T_{\text{interp}} + \big(\Sbar_{kgj} - \langle T_{\text{interp}}\rangle_{gj}\big).
\end{equation}
The pass then solves for the residual with $\That$ in place of $T$. Its bias is
$\Pi L^{-1}(\That - T)$, where $\That-T$ has zero mean on every bin and is the
interpolation error of the in-bin shape: second order in the node spacing, and much
smaller than the shape $T-\Sbar$ that the pass corrects (Sec.~6).

\section{Sampling}

The cell term $r_v = Q + \That - \Sigt(E)\tpsi$ is sampled as a single term. On a
sub-cell $c$ (energy interval $e$ of material $m$, a $\mu$ and a $\varphi$
interval of patch $j$),
\begin{equation}
  r_v = r_{\text{lin}} + \big(\Sigma_{\text{lin}}(E) - \Sigt(E)\big)\tpsi ,
\end{equation}
where $\Sigma_{\text{lin}}$ is linear between the interval's end nodes and
$r_{\text{lin}}$ is the trilinear interpolant of $r_v$'s node values. A trilinear
interpolant is a convex combination of the corner values, so
$|r_{\text{lin}}|\le\max_{\text{corners}}|r_v|$, and the proposal mass
\begin{equation}
  m_c = |c|\Big(\max_{\text{corners}}|r_v| + |\tpsi|\,\overline{|\Sigt-\Sigma_{\text{lin}}|}_e\Big),
\end{equation}
with the overbar the mean over the interval (exact on the cross-section grid), is
positive wherever $r_v$ can be nonzero. Within a sub-cell the point is uniform, and
$r_v$ is evaluated there with the pointwise $\Sigt(E)$. Sampling in proportion to
(a bound on) $|r|$ is the zero-variance choice for a single integral [11], applied
here to the source of a transport problem with signed weights [12]. The face term
keeps the proposal of \texttt{cartesian\_3d.tex}. The sub-cell masses and the face
masses share one CDF, so each term receives histories in proportion to its mass,
the Neyman allocation of \texttt{method/residual\_allocation.tex} with the volume mass
a bound rather than exact. A particle drawn from mass $m_n$ (total $m$) carries
$w = r\,|c|\,m/m_n$ before MC/DC's division by the number of histories.

\section{Verification}

On the $A=1$ infinite medium of \texttt{cartesian\_3d.tex} ($G=21$, octant patches,
$n_E=4$, $n_\mu=n_\varphi=3$):
\begin{itemize}
  \item \emph{Table.} The constant $\sigma_s$ and the isotropic centre-of-mass law
  give $\sum_{j'}U_{g'j'}(E) = N\sigma_s\ln(\min(E_{g'+1},\cdot)/\max(E_{g'},E))$
  exactly. The tabulated values agree to $1.6\times10^{-7}$ relative, and for an
  isotropic $\tpsi$, $T$ is the same at every direction node to $10^{-14}$.
  \item \emph{Interpolation.} Before the shift, the interpolant's bin averages
  differ from $\Sbar$ by about $10^{-4}$ relative ($2.5\times10^{-3}$ in the bin just
  below the source line, where $T$ falls to zero at the top edge), against an in-bin
  shape $T-\Sbar$ of order 10\%.
  \item \emph{Passes} (\texttt{test\_infinite\_medium\_A1}). The fixed point's error
  against the analytic bin averages runs from $-1.1\times10^{-3}$ to
  $+1.3\times10^{-3}$ across the bins. After eight passes of $3.4\times10^4$
  histories, the corrected averages agree with the analytic ones (largest $|z|=3.3$
  from the spread of the passes), and the coefficient of the fixed point's error
  pattern in the corrected error is $0.31\pm0.23$ (it is 1 without correction). One
  pass has a relative standard deviation of about $1.6\times10^{-3}$ per bin, and
  costs about one iteration; the table took 0.2\,s.
\end{itemize}
Fission enters $U$ through the isotropic-lab case. It is exercised by the fast
window of the fuel-pin problems, not by this test.

\section{Limitations}

\begin{itemize}
  \item The interpolation error is a deterministic bias. It is second order in the
  node spacing, and $n_E$, $n_\mu$, $n_\varphi$ (\texttt{correction\_nodes}) control
  it. Forward-peaked scattering (hydrogen) puts structure near the patch edges that
  needs more direction nodes than smooth kernels do.
  \item Memory: per material, (energy nodes) $\times$ (fundamental direction nodes)
  $\times\,G\times J$ values; about 46\,MB for $G=100$, $J=8$ and $3\times3$
  direction nodes. Evaluating $T$ at the nodes is one dense product per material and
  pass.
  \item The volume proposal uses the largest corner value, an upper bound on
  $|r_{\text{lin}}|$. Positive and negative parts of the in-bin shape are transported
  independently; pairing them (antithetic sampling) could reduce the variance
  further.
\end{itemize}

\begin{thebibliography}{99}
\bibitem{evans} T.~M.~Evans, T.~J.~Urbatsch, H.~Lichtenstein and J.~E.~Morel,
  ``A residual Monte Carlo method for discrete thermal radiative diffusion,''
  \emph{J.\ Comput.\ Phys.} \textbf{189}(2), 539--556 (2003).
\bibitem{peterson} J.~R.~Peterson, J.~E.~Morel and J.~C.~Ragusa,
  ``Exponentially-convergent Monte Carlo for the 1-D transport equation,'' M\&C 2013.
\bibitem{larsen} M.~Larsen et al., ``A Residual Monte Carlo Algorithm for Continuous
  Energy Neutron Transport with Elastic Scattering,'' \emph{Nucl.\ Sci.\ Eng.}
  \textbf{200}(3), 525--538 (2026), doi:10.1080/00295639.2025.2495608.
\bibitem{owen} A.~B.~Owen, \emph{Monte Carlo Theory, Methods and Examples} (2013).
\bibitem{bell} G.~I.~Bell and S.~Glasstone, \emph{Nuclear Reactor Theory}, Van
  Nostrand Reinhold (1970).
\bibitem{golub} G.~H.~Golub and J.~H.~Welsch, ``Calculation of Gauss quadrature
  rules,'' \emph{Math.\ Comp.} \textbf{23}(106), 221--230 (1969).
\bibitem{kronrod} A.~S.~Kronrod, \emph{Nodes and Weights of Quadrature Formulas},
  Consultants Bureau, New York (1965).
\bibitem{quadpack} R.~Piessens, E.~de Doncker-Kapenga, C.~W.~\"Uberhuber and
  D.~K.~Kahaner, \emph{QUADPACK: A Subroutine Package for Automatic Integration},
  Springer (1983).
\bibitem{duderstadt} J.~J.~Duderstadt and L.~J.~Hamilton, \emph{Nuclear Reactor
  Analysis}, Wiley (1976).
\bibitem{atkinson} K.~E.~Atkinson, \emph{An Introduction to Numerical Analysis}, 2nd
  ed., Wiley (1989).
\bibitem{kahn} H.~Kahn and A.~W.~Marshall, ``Methods of reducing sample size in Monte
  Carlo computations,'' \emph{J.\ Oper.\ Res.\ Soc.\ Am.} \textbf{1}(5), 263--278
  (1953).
\bibitem{lux} I.~Lux and L.~Koblinger, \emph{Monte Carlo Particle Transport Methods:
  Neutron and Photon Calculations}, CRC Press (1991).
\end{thebibliography}

\end{document}
